\SourceLeaf{196}\Needspace{10\baselineskip}\begin{center}
\textsc{Bemerkung.}
\end{center}

In einem Notizbuche von \textsc{Gauss} ist unter Aufzeichnungen aus den
Jahren 1824 bis 1828 vermerkt:
\begin{center}»\emph{Journal de l'École Normale. Fouriers
Defin. der Ebene. Lacroix I. p. 503.}«\end{center}

In der That enthalten die Séances des Écoles Normales T. I (Séance du 25
pluviose an III (1795)), Nouvelle Édition, Paris 1800, S. 28 Äusserungen
\textsc{Fouriers}, in denen er neue Erklärungen der Geraden und der Ebene
vorschlägt; \textsc{Fourier} beginnt dabei, ähnlich wie \textsc{Bolyai} und
\textsc{Lobatschefskij}, mit der Kugel. Weitere Veröffentlichungen
\textsc{Fouriers} über die Grundlagen der Geometrie scheinen nicht zu
existiren.

Der Discours préliminaire der Éléments de Géométrie von \textsc{Lacroix}
(Quatrième édition. Paris, An VIII [1804] S. \textsc{xv}.) enthält einen Hinweis auf
das \emph{Journal des Séances de l'École Normale}; wenn dabei nicht
\textsc{Fourier}, sondern \textsc{Laplace} genannt wird, so erklärt sich das
daraus, dass jener seine Theorie in den von diesem geleiteten geometrischen
Übungen der École Normale vorgetragen hatte.

\begin{flushright}
\textsc{Stäckel.}
\end{flushright}

\Needspace{10\baselineskip}
\medskip

\begin{center}
\textsc{Gauss} an \textsc{Schumacher}. Göttingen, 11. October 1827.
\end{center}

.\ .\ .\ .\ .\ .\ .\ .\ .\ Wer ist wohl Hr. \textsc{Koch} auf dem Cremon Nr. 83? Er hat
mir seine Parallelentheorie\HistoricalNote{[*)]}{[*)}{Christian Adolf \textsc{Koch}, \emph{Über Parallellinien.} Ein Versuch, dem Urtheil Sachkundiger gewidmet. Hamburg 1827. $8^{0}$. 12 S.]} zugeschickt, an der freilich nichts ist; aber
etwas seltenes ist, dass er meine Anzeige der Blösse sogleich anerkannt hat;
eben so wie meine Bemerkungen über seine nachher mir schriftlich
zugesandten, aber eben so erfolglosen Versuche. .\ .\ .\ .\ .

